\section{Generative Surrogate Models for Physical Fields}

Generative models aim to learn a probability distribution from reference samples and then generate new realizations from this distribution. Diffusion models are particularly relevant because they can represent complex high-dimensional distributions and have shown strong empirical performance in image generation. For physical surrogate modeling, the challenge is not only visual realism but also the preservation of statistical and physical properties.

Denoising Diffusion Probabilistic Models provide the main generative baseline
used in this report. Ho, Jain, and Abbeel introduced a formulation in which a
model learns to reverse a gradual noising process, making it possible to sample
complex high-dimensional distributions from noise \citep{Ho2020DDPM}. Song et
al. later connected diffusion and score-based generative modeling through
stochastic differential equations, giving a broader continuous-time view of the
same family of methods \citep{Song2021SDE}. The present work differs from these
general image-generation settings by applying diffusion models to RTI density
fields generated by DNS. The contribution is therefore not only to generate
visually plausible samples, but also to respect the physical symmetries of the
dataset through adapted augmentations, to evaluate generated fields with
physics-oriented metrics such as PhyFID and profile or spectral statistics, and
to explore wavelet conditioning as a scale-aware representation for turbulent
fields.

